% GENERATED FILE. Edit canonical lessons, then run npm run generate:books.
% canonical-source-hash: fnv1a64:0076d0f396ca573b
% canonical-lessons: TE-C159-pindu, TE-C159-nanabettu, TE-C159-baguceyu, TE-C159-pareyu, TE-C159-pogottuko

\chapter{To squeeze, To soak, To repair, To throw away, To lose}
\label{ch:te-to-squeeze-to-soak-to-repair-to-throw-away-to-lose}
\hlchaptermodality{159}

\begin{chapteropening}
\textbf{By the end of this chapter:} \emph{I can say to squeeze, to soak, to repair, to throw away and to lose.}

Complete the last lesson of chapter 159: I can say to squeeze, to soak, to repair, to throw away and to lose.
\end{chapteropening}

\section[\texorpdfstring{\te{పిండు} (piṇḍu)}{పిండు (piṇḍu)}]{\te{పిండు} (piṇḍu) — to squeeze}

\label{lesson:TE-C159-pindu}



\begin{quote}
Before the new one: say the Telugu for to take off, then the Telugu for to measure.
\end{quote}

\subsection*{You'll want to know: \te{పిండు}}
\textbf{\te{పిండు}} — \emph{piṇḍu} — \textquotedblleft{}to squeeze\textquotedblright{}.

\begin{grammarlens}[title={What you've built}]
One more word for this chapter.
\end{grammarlens}

\subsection*{Your turn}
\begin{itemize}
\raggedright
  \item \emph{Say it:} \emph{piṇḍu}
  \item \emph{Say it:} \emph{piṇḍu}, once more
  \item \emph{Say it:} say \emph{piṇḍu}
  \item \emph{Recall:} say the Telugu for to take off, then the Telugu for to measure, then say \emph{piṇḍu} again
\end{itemize}

\begin{tcolorbox}[breakable,colback=teal!4,colframe=teal!35!black,title={Before you move on}]
What does \te{పిండు} mean? (\textquotedblleft{}To squeeze\textquotedblright{}.) Say it once more.
\end{tcolorbox}
\section[\texorpdfstring{\te{నానబెట్టు} (nānabeṭṭu)}{నానబెట్టు (nānabeṭṭu)}]{\te{నానబెట్టు} (nānabeṭṭu) — to soak}

\label{lesson:TE-C159-nanabettu}



\begin{quote}
Before the new one: say the Telugu for to measure, then the Telugu for to squeeze.
\end{quote}

\subsection*{You'll want to know: \te{నానబెట్టు}}
\textbf{\te{నానబెట్టు}} — \emph{nānabeṭṭu} — \textquotedblleft{}to soak\textquotedblright{}.

\begin{grammarlens}[title={What you've built}]
One more word for this chapter.
\end{grammarlens}

\subsection*{Your turn}
\begin{itemize}
\raggedright
  \item \emph{Say it:} \emph{nānabeṭṭu}
  \item \emph{Say it:} \emph{nānabeṭṭu}, once more
  \item \emph{Say it:} say \emph{nānabeṭṭu}
  \item \emph{Recall:} say the Telugu for to measure, then the Telugu for to squeeze, then say \emph{nānabeṭṭu} again
\end{itemize}

\begin{tcolorbox}[breakable,colback=teal!4,colframe=teal!35!black,title={Before you move on}]
What does \te{నానబెట్టు} mean? (\textquotedblleft{}To soak\textquotedblright{}.) Say it once more.
\end{tcolorbox}
\section[\texorpdfstring{\te{బాగుచేయు} (bāgucēyu)}{బాగుచేయు (bāgucēyu)}]{\te{బాగుచేయు} (bāgucēyu) — to repair, to mend}

\label{lesson:TE-C159-baguceyu}



\begin{quote}
Before the new one: say the Telugu for to squeeze, then the Telugu for to soak.
\end{quote}

\subsection*{You'll want to know: \te{బాగుచేయు}}
\textbf{\te{బాగుచేయు}} — \emph{bāgucēyu} — \textquotedblleft{}to repair, to mend\textquotedblright{}.

\begin{grammarlens}[title={What you've built}]
One more word for this chapter.
\end{grammarlens}

\subsection*{Your turn}
\begin{itemize}
\raggedright
  \item \emph{Say it:} \emph{bāgucēyu}
  \item \emph{Say it:} \emph{bāgucēyu}, once more
  \item \emph{Say it:} say \emph{bāgucēyu}
  \item \emph{Recall:} say the Telugu for to squeeze, then the Telugu for to soak, then say \emph{bāgucēyu} again
\end{itemize}

\begin{tcolorbox}[breakable,colback=teal!4,colframe=teal!35!black,title={Before you move on}]
What does \te{బాగుచేయు} mean? (\textquotedblleft{}To repair, to mend\textquotedblright{}.) Say it once more.
\end{tcolorbox}
\section[\texorpdfstring{\te{పారేయు} (pārēyu)}{పారేయు (pārēyu)}]{\te{పారేయు} (pārēyu) — to throw away}

\label{lesson:TE-C159-pareyu}



\begin{quote}
Before the new one: say the Telugu for to soak, then the Telugu for to repair.
\end{quote}

\subsection*{You'll want to know: \te{పారేయు}}
\textbf{\te{పారేయు}} — \emph{pārēyu} — \textquotedblleft{}to throw away\textquotedblright{}.

\begin{grammarlens}[title={What you've built}]
One more word for this chapter.
\end{grammarlens}

\subsection*{Your turn}
\begin{itemize}
\raggedright
  \item \emph{Say it:} \emph{pārēyu}
  \item \emph{Say it:} \emph{pārēyu}, once more
  \item \emph{Say it:} say \emph{pārēyu}
  \item \emph{Recall:} say the Telugu for to soak, then the Telugu for to repair, then say \emph{pārēyu} again
\end{itemize}

\begin{tcolorbox}[breakable,colback=teal!4,colframe=teal!35!black,title={Before you move on}]
What does \te{పారేయు} mean? (\textquotedblleft{}To throw away\textquotedblright{}.) Say it once more.
\end{tcolorbox}
\section[\texorpdfstring{\te{పోగొట్టుకో} (pōgoṭṭukō)}{పోగొట్టుకో (pōgoṭṭukō)}]{\te{పోగొట్టుకో} (pōgoṭṭukō) — to lose (something)}

\label{lesson:TE-C159-pogottuko}



\begin{quote}
Before the new one: say the Telugu for to repair, then the Telugu for to throw away.
\end{quote}

\subsection*{You'll want to know: \te{పోగొట్టుకో}}
\textbf{\te{పోగొట్టుకో}} — \emph{pōgoṭṭukō} — \textquotedblleft{}to lose (something)\textquotedblright{}.

\begin{grammarlens}[title={What you've built}]
One more word for this chapter.
\end{grammarlens}

\subsection*{Your turn}
\begin{itemize}
\raggedright
  \item \emph{Say it:} \emph{pōgoṭṭukō}
  \item \emph{Say it:} \emph{pōgoṭṭukō}, once more
  \item \emph{Say it:} say \emph{pōgoṭṭukō}
  \item \emph{Recall:} say the Telugu for to repair, then the Telugu for to throw away, then say \emph{pōgoṭṭukō} again
\end{itemize}

\begin{tcolorbox}[breakable,colback=teal!4,colframe=teal!35!black,title={Before you move on}]
What does \te{పోగొట్టుకో} mean? (\textquotedblleft{}To lose (something)\textquotedblright{}.) Say it once more.
\end{tcolorbox}
